\documentclass[10pt,a4paper]{article}
\usepackage[margin=2cm]{geometry}
\usepackage{graphicx,booktabs,amsmath,amssymb,float,hyperref,caption}
\captionsetup{font=small}
\setlength{\parskip}{3pt}\setlength{\parindent}{0pt}
\newcommand{\fig}[3]{\begin{figure}[H]\centering\includegraphics[width=#2\linewidth]{results/figures/#1}\caption{#3}\end{figure}}
\title{ECM-402 Lab 7: Unsupervised Learning\\ICA, PCA and Clustering}
\author{Vivekanand Ojha \quad (Roll no.: \texttt{723159})}
\date{}
\begin{document}\maketitle

\textbf{Reproducibility.} Everything is regenerated by \texttt{python src/main.py}. Seeds: sources \texttt{seed(0)}; Gaussian sources \texttt{default\_rng(1)}; rotation baseline \texttt{default\_rng(123)}; FastICA \texttt{random\_state=0}; modulation dataset \texttt{seed=0}; KMeans/GMM/Spectral \texttt{random\_state=0} (KMeans \texttt{n\_init=10}); \texttt{make\_moons random\_state=0}. Covariances use $N-1$ throughout. The synthetic modulation dataset (not \texttt{load\_digits}) is used for Parts B--D.

\textbf{Label policy.} Ground-truth labels are used \emph{only} for evaluation (ARI, NMI, contingency table, matched correlation) or to colour plots. They are never given to any \texttt{fit()} and never used to choose $k$ or a number of components. The one exception, the supervised classifier in C12(c), is supervised by definition.

\section*{Part A --- ICA and blind source separation}
\subsection*{A1. Sources and mixtures}
\fig{A1_sources_mixtures.png}{0.8}{Sources (top) and mixtures (bottom).}
\textbf{(a)} Every mixture is a weighted sum of all three waveforms, so each $x_i$ shows the sine, square and sawtooth superimposed. The mixing weights are unknown and differ by channel, so no mixture looks like any single source and the sources cannot be read off by eye.

\textbf{(b)} $\mathrm{cond}(A)=9.85$ (singular values $3.167,\,0.982,\,0.321$; $\det A=1.0$). $A$ is full rank, so it is invertible, but a large condition number means the columns are nearly linearly dependent. Then $A^{-1}$ amplifies any noise or estimation error along the weak singular direction by up to $\mathrm{cond}(A)$, and the problem becomes ill-posed in the limit $\mathrm{cond}\to\infty$ (singular $A$: some sources are not recoverable at all). Here the conditioning is mild.

\subsection*{A2. Centering and whitening}
\textbf{(a)} Eigenvalues of $\mathrm{cov}(X)$: $0.107,\,0.915,\,10.242$. After $X_w=X_cE\Lambda^{-1/2}$, $\mathrm{cov}(X_w)=I_3$ to within $1.3\times10^{-14}$ (off-diagonals $\lesssim10^{-15}$; diagonals $1.0000000000000133,\,1.0000000000000004,\,1.0000000000000004$).

\textbf{(b)} Whitening removes all second-order structure (scales and correlations). In the ideal square, noiseless model, what remains is an orthogonal matrix, so ICA only has to search for a rotation: $n(n-1)/2=3$ free parameters for $n=3$ instead of $n^2=9$ for a general matrix (6 of the 9 are already fixed by whitening). This shrinks the search space and lets algorithms such as FastICA use fast orthogonal updates.

\subsection*{A3. FastICA}
\fig{A3_ica_recovered.png}{0.85}{True sources (left) and FastICA components (right, sign-aligned for display only).}
\textbf{(a, b)} Absolute-correlation matrix (rows $s_1,s_2,s_3$; columns IC0, IC1, IC2), with Hungarian one-to-one matching:
\[\begin{pmatrix}0.0001&\mathbf{0.9998}&0.0223\\0.0148&0.0182&\mathbf{0.9997}\\\mathbf{0.9996}&0.0295&0.0011\end{pmatrix}\]
Matches: $s_1\!\to$IC1 $0.99975$, $s_2\!\to$IC2 $0.99972$, $s_3\!\to$IC0 $0.99956$. The permutation and the signs differ from the original, as expected.

\textbf{(c)} Mean matched $|\rho|=\mathbf{0.99968}$. The remaining $\approx3\times10^{-4}$ gap is plausibly due to the added noise, the finite $N=2000$ and the fact that the periodic waveforms are only approximately independent; convergence took 7 iterations with no warnings. A stability check over FastICA seeds 0--9 gave $\approx0.9997$ for 8 seeds but $\approx0.837$ for seeds 4 and 5, i.e.\ one component was recovered poorly for those initialisations. I did not diagnose this further (probably a poor local solution), so a single FastICA run should not be trusted blindly.

\textbf{(d)} ICA determines each source only up to unknown scale and sign (and order). MSE would penalise a perfectly recovered but rescaled or sign-flipped source, whereas $|$Pearson correlation$|$ is invariant to both.

\subsection*{A4. PCA vs ICA}
\fig{A4_pca_vs_ica.png}{0.95}{$|$correlation$|$ with true sources for PCA and ICA, and off-diagonal $|D_{ij}|$.}
\textbf{(a)} PCA matched correlations: $0.620,\,0.671,\,0.816$, mean $\mathbf{0.702}$, versus $\mathbf{0.9997}$ for ICA. PCA only finds orthogonal directions of maximal variance, which depend on $A$'s scaling (here PC1 absorbs 10.2 of the total variance), and not the unmixing directions, so each PC remains a mixture of sources.

\textbf{(b)} Sample covariances (diagonals of the PCA output: $10.242,\,0.915,\,0.107$; ICA outputs: $1.0005$ each): both are diagonal to within $\sim10^{-15}$. So both output sets are uncorrelated, which says nothing about independence.

\textbf{(c)} $D_{ij}$ computed on unit-variance outputs (so scales are comparable):
\[D_{\text{PCA}}=\begin{pmatrix}0.908&-0.573&-0.265\\-0.573&1.710&-0.510\\-0.265&-0.510&1.522\end{pmatrix},\quad
D_{\text{ICA}}=\begin{pmatrix}0.811&0.0139&0.0080\\0.0139&0.523&0.0155\\0.0080&0.0155&0.0337\end{pmatrix}.\]
Mean $|D_{ij}|$ off-diagonal: PCA $0.450$, ICA $0.0124$ (the true sources themselves: $0.0065$). PCA outputs have large squared-component correlation, so they are clearly statistically dependent even though uncorrelated. ICA outputs are near the level of the true sources. A small value does not prove independence.

\textbf{(d)} PCA enforces only decorrelation (second order); independence requires the joint distribution to factorise, which constrains all orders. (b) shows PCA and ICA are equally uncorrelated; (c) shows only ICA is also free of this fourth-order dependence. FastICA does not minimise $D_{ij}$ directly; it maximises non-Gaussianity.

\subsection*{A5. Gaussian sources (non-identifiable)}
\fig{A5_gaussian_baseline.png}{0.55}{Gaussian case: random-rotation baseline vs FastICA scores.}
\textbf{(a)} Over 200 random rotations of the whitened Gaussian mixtures: mean matched $|\rho|=0.827$, 5th percentile $0.694$, 95th percentile $0.955$. Chance matching is far from 0 for $3\times3$ matrices.

\textbf{(b)} FastICA on $X_{\text{gauss}}$: mean matched $|\rho|=0.807$ (12 iterations, no convergence warning), inside the baseline range. The original illustrative signals give $0.99968$. Across FastICA seeds 0--9 the Gaussian score varied between $0.806$ and $0.884$, i.e.\ the answer depends on initialisation, which is the behaviour expected when no rotation is preferred. (Neither number was assumed to fall in a given range in advance.)

\textbf{(c)} For i.i.d.\ Gaussian sources any orthogonal rotation $Q$ of the whitened data is again jointly Gaussian with identity covariance, so $s$ and $Qs$ have exactly the same distribution; every rotation explains the observations equally well, and second-order statistics (all that whitened Gaussian data contains) cannot distinguish them. The sources are therefore identifiable only up to an arbitrary rotation, i.e.\ not identifiable, which is why ICA requires at most one Gaussian source.

\textbf{(d)} Excess kurtosis (\texttt{fisher=True}), sample (theory): original signals $-1.479\ (-1.5)$, $-1.968\ (-2.0)$, $-1.191\ (-1.2)$; Gaussian sources $-0.046,\,0.082,\,0.111$ (theory $0$). Strongly non-zero kurtosis marks non-Gaussian directions: a rotation mixes sources and pushes the kurtosis of each output toward 0, so maximising $|$kurtosis$|$ (or negentropy) singles out the true directions. Sample values deviate from theory for finite $N$ (here $N=2000$: the standard error of Gaussian sample kurtosis is $\approx\sqrt{24/N}\approx0.11$, matching the observed $\pm0.1$), and the added noise slightly shifts the non-Gaussian cases.

\section*{Part B --- PCA}
\subsection*{B6. Standardisation and dimensionality}
\textbf{(a)} Standardise first because PCA maximises variance, so features with larger numerical scale would dominate the components regardless of information content.

\fig{B6_scree.png}{0.9}{Scree plot and cumulative explained variance.}
\textbf{(b)} Explained variance ratios: $0.574,\,0.231,\,0.102,\,0.033,\,0.025,\,0.018,\,0.017$, then $\approx0$ for the rest. Cumulative: 80.5\% at $q=2$, 90.7\% at $q=3$, 94.0\% at $q=4$, 96.5\% at $q=5$. Minimum components: \textbf{3 for 90\%, 5 for 95\%}. Note the 12 features are far from independent: after unit-power normalisation $\mathrm{mean}|r|^2\equiv1$ (constant), and $|C_{40}|/C_{21}^2$, $|C_{42}|/C_{21}^2$ duplicate $|C_{40}|,|C_{42}|$ since $C_{21}=1$, so only about 7 directions carry variance.

\textbf{(c)} Not necessarily, and one should not expect it (here it is 3, not 4). The number of classes is the number of groups; intrinsic dimensionality is the number of directions along which the features vary. Classes can differ along a few shared axes (e.g.\ one axis separating BPSK from QAM), and within-class variation (noise, phase) also contributes variance unrelated to class identity.

\subsection*{B7. Reconstruction error}
\fig{B7_reconstruction.png}{0.55}{Reconstruction MSE vs $q$.}
\begin{center}\begin{tabular}{rrr}\toprule $q$ & measured MSE & $\frac{N-1}{N}\sum_{i>q}\lambda_i$\\\midrule
1&4.688053&4.688053\\2&2.144244&2.144244\\3&1.021504&1.021504\\5&0.384658&0.384658\\7&0.002401&0.002401\\\bottomrule\end{tabular}\end{center}
They agree to $\sim10^{-15}$ for all $q$ (full table in \texttt{results/tables/B7\_reconstruction\_mse.csv}). Reconstruction error is exactly the variance in the discarded directions, and since the eigenvectors capture the largest variances first, no other rank-$q$ linear projection can achieve a lower least-squares error (Eckart--Young).

\subsection*{B8. Two-component projection}
\fig{B8_pca_scatter.png}{0.5}{First two PCs. \textbf{Labels are used only for colouring; PCA was fitted without them.}}
\textbf{(a)} Done; labels only colour the plot.

\textbf{(b)} BPSK and 16QAM are well separated from everything else. QPSK and 8PSK are adjacent: centroid distance is only about 3.9 within-class spreads (vs.\ $\ge12.6$ for all other pairs), with a small overlap of the tails. Two PCs hold only 80.5\% of the variance, so discriminative information in the remaining 19.5\% is discarded, and low-variance directions can still be the discriminative ones.

\textbf{(c)} Largest $|$loadings$|$ on PC1: $|C_{42}|$ ($0.386$), $|C_{20}|$ ($0.352$), $|C_{40}|$ ($0.322$) (the normalised cumulant ratios duplicate these). These make physical sense: $|C_{20}|\approx1$ only for real-valued BPSK and $\approx0$ for the others; $|C_{40}|$ is about $2,1,0,0.68$ and $|C_{42}|$ about $2,1,1,0.68$ for BPSK/QPSK/8PSK/16QAM, so PC1 is largely a constellation-order axis (BPSK far right). The PSK vs QAM split is carried by PC2, whose largest loadings are the amplitude-spread features $\mathrm{std}|r|$, $\mathrm{std}|r|^2$ ($\approx+0.43$): PSK has constant envelope while 16-QAM has several amplitude levels.

\section*{Part C --- k-means}
\subsection*{C9. Clustering at $k=4$}
\textbf{(a)} Inertia $=2088.2$, silhouette $=0.480$. \textbf{(b)} ARI $=0.9973$, NMI $=0.9953$ (labels used only here for evaluation).
\textbf{(c)} Contingency table (rows true, columns cluster index):
\begin{center}\begin{tabular}{lrrrr}\toprule&c0&c1&c2&c3\\\midrule BPSK&0&249&0&1\\QPSK&0&0&0&250\\8PSK&250&0&0&0\\16QAM&0&0&250&0\\\bottomrule\end{tabular}\end{center}
Cluster indices are arbitrary labels assigned by initialisation, so plain accuracy would depend on a meaningless permutation unless clusters are first optimally matched to classes.

\textbf{(d)} The Rand index is the fraction of sample pairs on which two partitions agree; ARI subtracts the value expected for random partitions with the same cluster sizes and rescales, so $1$ is perfect agreement and $\approx0$ is chance. For a single random assignment ARI is a random variable with mean $\approx0$ but nonzero spread (empirically, mean $-5\times10^{-5}$, s.d.\ $0.0013$ for 200 random 4-way assignments here), so one draw need not be exactly 0.

\subsection*{C10. Choosing $k$ without labels}
\fig{C10_k_selection.png}{0.95}{Inertia, silhouette and (evaluation-only) ARI vs $k$.}
\begin{center}\begin{tabular}{rrrr}\toprule $k$&inertia&silhouette&ARI (eval only)\\\midrule
2&5581.4&0.509&0.333\\3&2531.9&\textbf{0.588}&0.714\\4&2088.2&0.480&\textbf{0.997}\\5&1868.4&0.467&0.877\\6&1672.0&0.399&0.786\\7&1505.7&0.390&0.711\\8&1352.7&0.394&0.699\\\bottomrule\end{tabular}\end{center}
\textbf{(b)} Using internal criteria only I would pick \textbf{$k=3$}: the silhouette is highest there ($0.588$), and the elbow (largest second difference of the inertia) is also at 3, though the bend at 4 is weaker but visible, so the elbow plot alone is not decisive.
\textbf{(c)} Highest silhouette is at $k=3$; highest ARI at $k=4$. They disagree. Inspecting the $k=3$ partition shows it merges QPSK and 8PSK into one cluster (BPSK and 16QAM stay separate), exactly the close pair seen in B8; merging gives compact, well-separated clusters, so silhouette rewards it although it loses a true class.
\textbf{(d)} Without labels, a single internal index can be systematically wrong when classes are close. A practical safeguard: combine several indices (e.g.\ silhouette, Davies--Bouldin, Calinski--Harabasz, gap statistic), check stability across resamples/seeds, and inspect the clusters with domain knowledge.

\subsection*{C11. PCA before clustering}
\begin{center}\begin{tabular}{lrrr}\toprule space&ARI&NMI&silhouette\\\midrule 12-D standardised&0.9973&0.9953&0.480\\2-D PCA&0.9867&0.9805&0.744\\\bottomrule\end{tabular}\end{center}
Reducing to 2-D did not improve ARI (it dropped slightly, $0.997\to0.987$, i.e.\ a handful of samples: the 2-D PCA discards 19.5\% of variance that helps separate QPSK/8PSK). The silhouette jumps because distances are measured in a different space: the 12-D space has redundant and noisy dimensions that add within-cluster spread. Scoring the \emph{12-D clustering} in the 2-D space already gives $0.742$, so almost all of the gain is from the space, not the clustering. Silhouette values from different feature spaces are not comparable.

\subsection*{C12. ARI vs SNR}
\fig{C12_snr_sweep.png}{0.9}{k-means ARI vs SNR (seed-0 dataset) and a supervised stand-in.}
\begin{center}\begin{tabular}{rrrrr}\toprule SNR (dB)&ARI (seed 0)&ARI mean$\pm$sd (10 datasets)&NMI (seed 0)&Supervised acc.\\\midrule
20&0.997&$0.9997\pm0.0008$&0.995&1.00\\15&1.000&$0.967\pm0.103$&1.000&1.00\\10&1.000&$0.799\pm0.173$&1.000&1.00\\5&0.630&$0.620\pm0.009$&0.771&0.99\\0&0.346&$0.321\pm0.018$&0.409&0.65\\\bottomrule\end{tabular}\end{center}
\textbf{(b)} The single seed-0 curve is flat to 10 dB and then drops sharply, but that is partly luck: averaged over 10 independently generated datasets the degradation starts earlier (15 dB: $0.97$, 10 dB: $0.80\pm0.17$) and is large at 5 dB ($0.62$). Random clustering gives ARI $\approx0$ (here $-5\times10^{-5}\pm0.0013$); at 0 dB ARI is $0.32$--$0.35$, so still clearly above chance.
\textbf{(c)} I do not have the Lab 3 curve in this repo, so the supervised column is my own stand-in (standardised logistic regression on the same features, 70/30 stratified split, seed 0); replace it with your Lab 3 accuracies if they differ. Using this stand-in, clustering falls behind from about 10 dB (dataset-averaged) / 5 dB (seed 0), where supervised accuracy is still $\ge0.99$. A supervised model can place boundaries between classes whose features are shifted by noise but still consistently different, even when clusters overlap, whereas k-means needs compact, well separated groups to exist in the first place. At 0 dB even the supervised stand-in drops to $0.65$.

\section*{Part D --- When k-means is the wrong tool}
\subsection*{D13. Algorithm comparison (20 dB, $k=4$, standardised 12-D)}
Seeds: \texttt{KMeans(n\_init=10, random\_state=0)}, \texttt{GaussianMixture(covariance\_type="full", random\_state=0)} (default \texttt{n\_init=1}, \texttt{reg\_covar=1e-6}), \texttt{AgglomerativeClustering(linkage="ward")} (deterministic).
\begin{center}\begin{tabular}{lrrrr}\toprule method&ARI&NMI&ARI over 10 datasets&NMI over 10 datasets\\\midrule
KMeans&0.9973&0.9953&$0.9997\pm0.0008$&$0.9995\pm0.0015$\\GMM (full)&1.0000&1.0000&$0.905\pm0.152$&$0.932\pm0.109$\\Agglomerative (ward)&1.0000&1.0000&$0.978\pm0.048$&$0.982\pm0.039$\\\bottomrule\end{tabular}\end{center}
\textbf{(a)} In the one required run, GMM and Ward both reach 1.0 and k-means $0.9973$, but this gap is one sample (one BPSK point) and does not support any ranking. A single run reflects one dataset draw and one initialisation. The supplementary 10-dataset experiment (seeds 0--9; the dataset seed is also used as the model seed) reverses the ranking: k-means is best and very stable, Ward is next, and the full-covariance GMM is worst and most variable. A likely (untested) reason is that the full covariance of GMM is estimated on features that are nearly collinear or duplicated (see B6), making some component covariances almost singular, plus a single initialisation.
\textbf{(b)} A GMM gives each cluster its own full covariance matrix, so clusters can be elliptical with different orientations and sizes, and it uses soft assignments with mixture weights instead of hard nearest-centroid assignment.

\subsection*{D14. Two moons}
\fig{D14_moons.png}{0.8}{k-means vs spectral clustering on \texttt{make\_moons(500, noise=0.05)}.}
\textbf{(a)} k-means ($k=2$) gets ARI $=0.257$ (NMI $0.195$): it cuts the data with a straight line. A centroid-based rule assigns each point to the nearest centroid, so the decision boundary is the perpendicular bisector of the two centroids, a straight line; the two interleaved crescents are non-convex, and no straight line can separate them, since each crescent contains points on both sides of any such line through the middle.
\textbf{(b)} Spectral clustering (nearest-neighbour affinity) gives ARI $=1.0$ (NMI $1.0$). It builds a neighbourhood graph in which points are linked to nearby points and clusters along the crescents are connected, then splits the graph using eigenvectors of the graph Laplacian, so cluster membership follows connectivity rather than distance to a centre. (scikit-learn warned that the graph was not fully connected; the result was nevertheless perfect here.)

\end{document}
